\chapter{}
\label{chap:1}

\section{Умова}
Оцінити часову складність та обсяг матеріалу, необхідного для проведення атаки
Куртуа-Майєра на фільтрувальний генератор гами, який складається з ЛРЗ довжини $n$
та функції ускладення:
\begin{equation*}
    f(x_{1}, x_{2}, x_{3}, x_{4}) = x_{1}x_{2}x_{3} \oplus x_{1}x_{2} \oplus x_{1}x_{3} \oplus x_{1} \oplus x_{2}.
\end{equation*}

\section{Розв'язання}
Схожого типу задача була в дз \No 4, то я використосував вже існуючі таблиці, просто змінюючи значення.

\noindent $D = \deg f = 3$, $n = 4$. Множина нулів та одиниць для даної функції $f$ має вигляд:
\begin{center}
    \begin{tabular}{c|cccccccc}
        $(x_1, x_2, x_3, x_4)$ & $0000$ & $0001$ & $0010$ & $0011$
                               & $0100$ & $0101$ & $0110$ & $0111$                         \\
        \hline
        $f$                    & $0$    & $0$    & $0$    & $0$    & $1$ & $1$ & $1$ & $1$
    \end{tabular}
\end{center}
\begin{center}
    \begin{tabular}{c|cccccccc}
        $(x_1, x_2, x_3, x_4)$ & $1000$ & $1001$ & $1010$ & $1011$
                               & $1100$ & $1101$ & $1110$ & $1111$                         \\
        \hline
        $f$                    & $1$    & $1$    & $0$    & $0$    & $1$ & $1$ & $1$ & $1$
    \end{tabular}
\end{center}
\begin{align*}
    & V (I) = \{(0,0,0,0),\; (0,0,0,1),\; (0,0,1,0),\; (0,0,1,1),\; (1,0,1,0),\; (1,0,1,1)\}, \\
    & V_{n} \setminus V(I) = \{(0,1,0,0),\; (0,1,0,1),\; (0,1,1,0),\; (0,1,1,1),\; (1,0,0,0), \\
    & (1,0,0,1),\; (1,1,0,0),\; (1,1,0,1),\; (1,1,1,0),\; (1,1,1,1)\}.
\end{align*}
Отже: $\operatorname{wt}(f) = 10$, $\operatorname{wt}(f \oplus 1) = 6$.

Щоб обчислити алгебраїчну імунність, треба порахувати дві величини. Маємо $\AI(f) \leq \lceil \frac{n}{2} \rceil = 2$, 
а тому перевіряємо, чи існують ненульові функції меншого степеня (в даному випадку $1$) в $\Ann(f)$.

Нехай $g = a_0 \oplus a_1 x_1 \oplus a_2 x_2 \oplus a_3 x_3 \oplus a_4 x_4 \in \Ann(f)$ -- в загальному випадку, тобто $g$
має обертатися в нуль на $V_{n} \setminus V(I)$. Підставляємо:
\begin{equation*}
    \begin{cases}
        g(0,1,0,0) = a_0 \oplus a_2 = 0,                                    \\
        g(0,1,0,1) = a_0 \oplus a_2 \oplus a_4 = 0 \;\Rightarrow\; a_4 = 0, \\
        g(0,1,1,0) = a_0 \oplus a_2 \oplus a_3 = 0 \;\Rightarrow\; a_3 = 0, \\
        g(1,0,0,0) = a_0 \oplus a_1 = 0,                                    \\
        g(1,1,0,0) = a_0 \oplus a_1 \oplus a_2 = 0.                         \\
    \end{cases}
\end{equation*}
Ксоримо 1 та 2 рівняння: $a_4 = 0$. Проксоривши 1 та 3: $a_3 = 0$. Проксоривши 1 та 5: $a_1 = 0$. З 4-го: 
$a_1 = a_0 = 0$. Підставимо в перше: $a_2 = a_0 = 0$. Маємо єдиний розв'язок, який є тривіальним, тобто 
ненульових ануляторів степеня $1$ в $\Ann(f)$ не існує. Оскільки $f \oplus 1 \in \Ann(f)$ і $\deg(f \oplus 1) = 3$, 
то маємо $\min \deg \Ann(f) \geq 2$.

Далі аналогічним чином обчислюємо $\min \deg \Ann(f \oplus 1)$. Покладемо $h$, що має вигляд 
$h = a_0 \oplus a_1 x_1 \oplus a_2 x_2 \oplus a_3 x_3 \oplus a_4 x_4 \in \Ann(f \oplus 1)$, тобто $h$ обертається 
в нуль на $V(I)$:
\begin{equation*}
    \begin{cases}
        h(0,0,0,0) = a_0 = 0,                                                 \\
        h(0,0,0,1) = a_0 \oplus a_4 = 0 \;\Rightarrow\; a_4 = 0,              \\
        h(0,0,1,0) = a_0 \oplus a_3 = 0 \;\Rightarrow\; a_3 = 0,              \\
        h(0,0,1,1) = a_0 \oplus a_3 \oplus a_4 = 0,                           \\
        h(1,0,1,0) = a_0 \oplus a_1 \oplus a_3 = 0 \;\Rightarrow\; a_1 = 0,   \\
        h(1,0,1,1) = a_0 \oplus a_1 \oplus a_3 \oplus a_4 = 0.
    \end{cases}
\end{equation*}
З першого рівняння одразу $a_0 = 0$, далі послідовно $a_4 = 0$, $a_3 = 0$, $a_1 = 0$. Коефіцієнт $a_2$ виходить 
є вільним, бо жодне рівняння його не обмежує. То маємо два випадки -- тривіальний відкидаємо поки що, і покладаємо 
$a_2 = 1$, тоді $h = x_2$, $\deg h = 1$. Перевіримо, чи обертається це в нуль на $V(I)$. Всі вектори у множині 
$V (I) = \{(0,0,0,0),\; (0,0,0,1),\; (0,0,1,0),\; (0,0,1,1),\; (1,0,1,0),\; (1,0,1,1)\}$ мають $x_2 = 0$, а тому 
$h = x_2 = 0$ на них. Отже, $h = x_2 \in \Ann(f \oplus 1)$, $\deg h = 1$, тобто $\min \deg \Ann(f \oplus 1) = 1$.

Тоді алгебраїчна імунність має бути мінімумом з двох значень:
\begin{equation*}
    \AI(f) = \min\big\{\min \deg \Ann(f),\; \min \deg \Ann(f \oplus 1)\big\} = \min\{\geq 2,\; 1\} = 1.
\end{equation*}
Для оцінки часової складності, при $d = \AI(f) = 1$, $D = \deg f = 3$, маємо, що:
\begin{equation*}
    N_{d} = N_{1} = \binom{n}{1} = n.
\end{equation*}
Необхідна кількість бітів гами (обсяг матеріалу): $t \geq N_1 = n$. А часова складність атаки складає:
\begin{equation*}
    T \approx O\!\left(t \cdot N_1^2\right) = O\!\left(N_1^3\right) = O(n^3).
\end{equation*}
Оскльки у нас $d < D$, то порівняно зі стандартною атакою (без пошуку ануляторів) має бути виграш. Мені аж стало 
цікаво який (бо в дз виграшу не було). Якби нам довелося розв'язувати вхідну систему степеня $D = 3$ методом 
введення нових змінних, кількість невідомих склала б:
\begin{equation*}
    N_D = N_3 = \binom{n}{1} + \binom{n}{2} + \binom{n}{3}
    = n + \frac{n(n-1)}{2} + \frac{n(n-1)(n-2)}{6}.
\end{equation*}
Тоді виграш у часовій складності визначиться співвідношенням:
\begin{equation*}
    \left(\frac{N_D}{N_d}\right)^{\!3} = \left(\frac{N_3}{N_1}\right)^{\!3}
    = \left(1 + \frac{n-1}{2} + \frac{(n-1)(n-2)}{6}\right)^{\!3} \approx n^{6}.
\end{equation*}
В даному випадку атака Куртуа-Майєра дає ну дуже суттєвий виграш. Це все через те, що існує анулятор 
$h = x_2 \in \Ann(f \oplus 1)$ степеня $1$, що значно менше від $\deg f = 3$.